\subsection{Research Gap: Positioning This Work}

The eighteen studies agree that a deep reinforcement learning agent can learn a profitable trading policy. Where they differ, or stay silent, is on the conditions under which that holds. The findings above describe what the field does; read against one another they also show what it has not yet done. Table~\ref{Tables:Positioning} sets out the comparison.

\input{Tables/Positioning}

The first gap appears only when two of those columns are read together. The three studies that trade currencies \cite{carapuco_reinforcement_2018, huang_financial_2018, rundo_deep_2019} all restrict the agent to a discrete choice, and the three that give the agent a continuous action \cite{majidi_algorithmic_2022, jia_lstm-ddpg_2021, zhang_deep_2020} all trade something else. Among the work reviewed here the two properties never occur together, so there is no direct comparison available for a currency agent that sets its own position size. That is the combination this work uses, and the absence is the reason it cannot be measured against an existing result.

The second gap follows from the cost figures. A strategy evaluated without charges is not evaluated against the market it claims to trade, and the omission matters more in Forex than the reported returns suggest, because the spread is not fixed. It widens when liquidity falls and around economic releases, which are the moments at which a model trained on an average spread would expect to act.

The third gap is resolution rather than length of history. A long sample is already the norm in this literature, so it is not by itself distinctive; working inside the bar is. A bar reports four prices and conceals the order in which the price visited them, and that order decides which spread was paid and whether a stop was reached before a target.

The fourth gap is one of evidence rather than of design. Each study reports the performance of a model it selected, and none of those reviewed here states how many candidates that model was selected from, or how the reported figure moves when the evaluation is repeated from a different starting condition. Deep reinforcement learning is sensitive to initialisation, so a single figure does not separate the contribution of the method from the contribution of the run, and a model chosen as the best of many describes the sample it was chosen from as much as the market it traded \cite{bailey_pseudo-mathematics_2014, lopez_de_prado_advances_2018}.

The last gap concerns how results are combined. The studies reviewed report either a single agent trading one instrument, or a single agent trading a basket. None of them trains one design separately on each of several instruments and then measures what those agents do when they are run together, which is how a trader holding several currency pairs would deploy them. The combination is not implied by the individual results, because it depends on how the equity curves move relative to one another.

This work is positioned in those gaps. One design is applied unchanged to all seven major currency pairs, with a continuous action that sets the direction and the size of the position together. The evaluation runs on quote-level data across eleven years and charges the spread quoted at the moment of each trade together with the commission of a swap-free raw-spread retail account, which is an account a retail trader can open rather than a modelling convenience. Each pair is reported from five different starting balances rather than from a single backtest, and the number of candidates the selected model was drawn from is stated. The seven agents are then combined into an equal-weight portfolio, so the work reports both the seven individual results and the single result a trader running all of them would have seen. Section~\ref{Section:ExperimentalResults} gives both.

One point of similarity is worth stating alongside the differences. The reward is the log return, which is the choice made by one of the surveyed studies \cite{majidi_algorithmic_2022} and is not a point of difference. What differs is that the reward the agent optimises and the criteria the final model is selected on are deliberately not the same measure. Candidate runs are scored during training by the Calmar ratio, and the model finally published for each pair is the one ranked highest by a composite of percentile ranks within its own pool of candidates, built from the annualised return, the Sharpe, Sortino, Calmar and Sterling ratios defined in Section~\ref{Subsection:PerformanceMeasurement}, the maximum and downside risk, and measures of how the policy behaves. A policy that earns its return by taking more risk therefore does not win on return alone. Section~\ref{Section:Implementation} sets out both.